% Section 1.4, from lectures/L01/appendix/c_exercises.md.
%
% One deliberate difference: exercise 10 a) spells the text IC74 rather than the name of the class
% the course was first given to, so that the book belongs to no one year.

\section{Övningar}
\label{c1:sec:exercises}\label{c1:app:c}

\begin{aside}
\textbf{Så kontrollerar du ditt arbete.} Alla övningar här räknas med papper och penna. Räkna varje
omvandling för hand först, och kontrollera den sedan genom att räkna tillbaka åt andra hållet: ett
binärt tal du har räknat fram ur ett decimalt ska ge samma decimala tal när du summerar vikterna.
Den kontrollen hittar nästan alla fel, och den finns alltid till hands, också på provet.

Lösningsförslagen finns i bilaga~\ref{app:answers}. Där visas också de vanligaste felaktiga
svaren, och varför de är fel.

Varje övning är märkt med sin sort. Övning~\ref{c1:ex:13} är kapitlets
\textbf{kontrolluppgift}: där räknar du för hand och kontrollerar med kalkylatorn i Windows.
\end{aside}

\exercise{Positionssystem}{Förståelse}
\label{c1:ex:1}
\begin{parts}
  \item Talet 4072 är decimalt. Ange för varje siffra dess position, dess vikt och det värde den
    bidrar med.
  \item I ett talsystem med talbasen 5 finns bara siffrorna 0--4. Vilken vikt har position 2 där?
    Vilket decimalt tal är $243_5$?
  \item Förklara med egna ord varför den högra positionen har vikten 1 oavsett talbas.
\end{parts}

\exercise{Från binärt till decimalt}{Räkna för hand}
\label{c1:ex:2}
Omvandla till decimalt. Använd viktraden 128, 64, 32, 16, 8, 4, 2, 1.
\begin{parts}
  \item \code{1101}
  \item \code{1 0000}
  \item \code{0110 0100}
  \item \code{1000 0001}
  \item \code{1111 1111}
  \item \code{1101 0110}
\end{parts}

\exercise{Från decimalt till binärt}{Räkna för hand}
\label{c1:ex:3}
Omvandla till binärt. Använd vikttabellen för a)--c) och division med 2 för d)--f), och
kontrollera varje svar genom att räkna tillbaka.
\begin{parts}
  \item 13
  \item 37
  \item 100
  \item 200
  \item 255
  \item 1000
\end{parts}

\exercise{Binärt och hexadecimalt}{Räkna för hand}
\label{c1:ex:4}
Omvandla från binärt till hexadecimalt:
\begin{parts}
  \item \code{1010 1111}
  \item \code{0011 1100}
  \item \code{1 1110}
  \item \code{11 1110 1000}
\end{parts}
Omvandla från hexadecimalt till binärt:
\begin{parts}[start=5]
  \item \code{0x7F}
  \item \code{0xC3}
  \item \code{0x2A5}
\end{parts}

\exercise{Hexadecimalt och decimalt}{Räkna för hand}
\label{c1:ex:5}
Omvandla från hexadecimalt till decimalt:
\begin{parts}
  \item \code{0x2A}
  \item \code{0xFF}
  \item \code{0x100}
  \item \code{0x3E8}
\end{parts}
Omvandla från decimalt till hexadecimalt, med valfri metod:
\begin{parts}[start=5]
  \item 75
  \item 160
  \item 4095
\end{parts}

\exercise{Bitar och talområden}{Förståelse}
\label{c1:ex:6}
\begin{parts}
  \item Hur många olika värden kan ett tal med 4, 8, 12 respektive 16 bitar anta, och vilket är det
    största?
  \item Hur många bitar behövs minst för att kunna representera 1000 olika värden?
  \item En nivågivare i en tank ska rapportera nivån i hela steg från 0 till 200. Hur många bitar
    behövs? Hur många av kombinationerna blir då oanvända?
  \item Utan att räkna ut värdet: är \code{1011 0111} ett jämnt eller ett udda tal? Är det större
    eller mindre än 128? Motivera.
\end{parts}

\exercise{Binär addition}{Räkna för hand}
\label{c1:ex:7}
Addera de 8-bitars binära talen. Skriv ut minnessiffrorna, och kontrollera varje summa genom att
göra om talen till decimalt.
\begin{parts}
  \item \code{0011 1010 + 0001 0110}
  \item \code{0111 1111 + 0000 0001}
  \item \code{1100 1000 + 0110 0100}
  \item Vilken av summorna i a)--c) ryms inte i åtta bitar? Vad blir kvar i ett 8-bitars register,
    och vad hände med resten?
\end{parts}

\exercise{BCD}{Räkna för hand}
\label{c1:ex:8}
\begin{parts}
  \item Koda talet 47 i BCD. Koda det sedan som ett 8-bitars binärt tal, och jämför.
  \item Koda årtalet 2026 i BCD. Hur många bitar behövs?
  \item Vilket decimalt tal är \code{1001 0011} i BCD? Vilket tal är samma bitmönster som binärt
    tal?
  \item Är \code{0110 1010} ett giltigt BCD-tal? Motivera.
\end{parts}

\exercise{Graykod}{Förståelse}
\label{c1:ex:9}
\begin{parts}
  \item Bygg Graykoden med tre bitar genom att spegla tvåbitarskoden \code{00, 01, 11, 10}, enligt
    \secref{c1:b3}.
  \item Hur många bitar ändras när en binär räknare går från 7 till 0 (\code{111} till
    \code{000})? Hur många ändras mellan motsvarande Graykoder?
  \item En vinkelgivare med tre bitar står mellan läge 1 och läge 2. Förklara vad styrsystemet kan
    läsa under övergången om givaren använder vanlig binärkod, och vad det kan läsa om den
    använder Graykod.
\end{parts}

\exercise{ASCII}{Räkna för hand}
\label{c1:ex:10}
Använd tabellen i \secref{c1:b4}.
\begin{parts}
  \item Skriv texten \code{IC74} som ASCII-koder, hexadecimalt.
  \item Vilken text är ASCII-koderna \code{0x48 0x65 0x6A}? (\code{e} ligger på \code{0x65}, och
    bokstäverna ligger i ordning.)
  \item Hur mycket skiljer \code{a} och \code{A}? Vilken enda bit skiljer dem åt?
  \item Ett program har läst in siffertecknet \code{7} från ett tangentbord. Hur får programmet
    fram talet 7 ur det? Och hur gör det om talet 3 till tecknet \code{3}?
\end{parts}

\exercise{7-segmentkod}{Konstruktion}
\label{c1:ex:11}
Använd bitordningen g f e d c b a från \secref{c1:b5}.
\begin{parts}
  \item Vilka segment ska lysa för den hexadecimala siffran \code{E}? Ange koden binärt och
    hexadecimalt.
  \item En display får koden \code{0x6D}. Vilken siffra visas?
  \item Segment g har gått sönder och lyser aldrig. Vilka siffror går då inte längre att skilja från
    någon annan siffra? Vilka ser bara konstiga ut?
\end{parts}

\exercise{Paritetsbit}{Räkna för hand}
\label{c1:ex:12}
Använd jämn paritet: paritetsbiten väljs så att det totala antalet ettor blir jämnt, och den skickas
som bit~7 framför de sju ASCII-bitarna.
\begin{parts}
  \item Ange den byte som skickas för tecknen \code{B}, \code{C} och \code{G}.
  \item En mottagare får byten \code{1100 0001}. Har ett fel inträffat? Kan mottagaren veta vilket
    tecken som skulle ha skickats?
  \item Visa med ett eget exempel att två felaktiga bitar i samma byte inte upptäcks.
\end{parts}

\exercise{Kontroll: för hand, och sedan med kalkylatorn}{Kontroll}
\label{c1:ex:13}
\emph{Räkna för hand, kontrollera med ett verktyg, och förklara det du ser.}

Gör delarna i ordning, och skriv ned varje svar innan du går vidare. Poängen går förlorad om du
tittar i kalkylatorn först.

\task{a) För hand}
Omvandla 173 till binärt och hexadecimalt, \code{0x3E8} till decimalt och binärt, och
\code{0b1100 1010} till decimalt och hexadecimalt.

\task{b) Med kalkylatorn}
Öppna kalkylatorn i Windows och välj läget \textbf{Programmerare}. Skriv in vart och ett av talen i
det talsystem det är givet i (klicka först på \code{DEC}, \code{HEX} eller \code{BIN}), och läs av
de andra två. Stämmer dina svar? Om inte: hitta felet i din uträkning, inte bara rätt svar.

\task{c) Något du inte kan förklara än}
Välj \code{DEC}, skriv \code{-1} (tryck \code{1} och sedan \code{+/-}), och läs av \code{HEX} och
\code{BIN}. Byt sedan ordlängd med knappen där det står \code{QWORD} tills det står \code{BYTE}, och
läs av igen. Skriv ned vad du ser i båda fallen.

\task{d)}
Vilket positivt tal har samma bitmönster som det kalkylatorn visade för \code{-1} i läget
\code{BYTE}? Gissa varför datorn skriver \code{-1} just så. Svaret, som heter
\textbf{tvåkomplement}, kommer i kapitel~\ref{c8:ch}.

\exercise{Att lägga till en nolla}{Förståelse, Fördjupning}
\label{c1:ex:14}
\begin{parts}
  \item \code{0b1011} är 11. Vilket tal är \code{0b10110}, där en nolla har lagts till längst till
    höger? Och \code{0b101100}?
  \item Formulera en regel: vad händer med värdet på ett binärt tal när alla bitar flyttas ett steg
    åt vänster och en nolla fylls på från höger?
  \item Vad händer i stället med värdet om den högra biten tas bort och alla bitar flyttas ett steg
    åt höger? Pröva med \code{0b1011} och \code{0b1010}.
  \item Vad motsvarar det att lägga till en nolla längst till höger i ett hexadecimalt tal? Pröva
    med \code{0x2A}.
\end{parts}
Mikrodatorn i boken har instruktioner som gör exakt det här med ett register, och du kommer att
använda dem i kapitel~\ref{c8:ch}.
